\label{sec:bisect-integration}

\subsection{The Three-Layer Compositor}

The \textsc{bisect} redistricting system structures every plan by three independently
configurable components:
\begin{enumerate}
  \item \textbf{Structure} (Layer 1): the bisection tree shape, which determines
    how the state is recursively divided into the required number of districts.
    Options include \texttt{standard-bisect}, \texttt{prime-factor} (ApportionRegions),
    and \texttt{ratio-optimal} (GeoSection).
  \item \textbf{Weights} (Layer 2): the edge weight mode, which determines how
    costly each potential cut is for the METIS partitioner.
    Current modes are \texttt{geographic} (boundary length), \texttt{county}
    (county co-membership bonus), \texttt{unweighted}, and \texttt{vra-aligned}.
  \item \textbf{Search} (Layer 3): the seed selection strategy, which determines
    which METIS result to accept when multiple random seeds are attempted.
    Options include \texttt{single}, \texttt{multi}, \texttt{convergence},
    and \texttt{percentile}.
\end{enumerate}
Community character weights occupy Layer 2.
They do not alter the bisection tree structure (Layer 1) or the seed selection
strategy (Layer 3).

\subsection{Existing Weight Modes}

The current weight modes are:
\begin{itemize}
  \item \textbf{Geographic}: $w(u,v)$ proportional to the shared boundary length
    between tracts $u$ and $v$.
    Long shared boundaries are more costly to cut; this encourages compact, round
    districts.
  \item \textbf{County}: a bonus weight applied to edges between tracts in the same
    county, as introduced in paper B.10.
    This encourages county-level community preservation without explicitly encoding
    any economic or cultural signal.
  \item \textbf{Unweighted}: $w(u,v) = 1$ for all edges.
    Used as a baseline; METIS minimizes edge count (number of tract boundary pairs
    split) rather than weighted cut cost.
  \item \textbf{VRA-aligned}: a weight mode that boosts edges between tracts with
    high minority population concentrations to preserve VRA Section 2 majority-minority
    configurations.
    See paper B.14 for details.
\end{itemize}

\subsection{Alpha-Blending Community Character with Base Weights}

The M-track introduces community character weights as an additional modifier that
composes with any base weight mode via \emph{alpha-blending}:
\begin{equation}
  w_{\mathrm{final}}(u,v)
  \;=\;
  \alpha \cdot w_{\mathrm{base}}(u,v)
  \;+\;
  (1 - \alpha) \cdot w_{\mathrm{char}}(u,v),
  \quad \alpha \in [0,1],
  \label{eq:alpha-blend}
\end{equation}
where $w_{\mathrm{base}}(u,v)$ is the weight from the active base mode
(geographic, county, or unweighted) and
$w_{\mathrm{char}}(u,v) = w_{\mathrm{base}}(u,v) \times \mathrm{sim}(\mathbf{c}_u, \mathbf{c}_v)$
is the community character modified weight.

When $\alpha = 1$, the final weight equals the base weight (community character
is inactive).
When $\alpha = 0$, the final weight is entirely determined by community character
similarity.
The default in M-track configurations is $\alpha = 0.5$, which gives equal weight
to geographic boundaries and community character.

\textbf{Operational formula:} Equation~\ref{eq:alpha-blend} is the formula
implemented in \texttt{crates/bisect-cli/src/edge\_weights.rs}
(\texttt{EconomicCharacterWeighter}). Equation~\ref{eq:general-similarity} in
\S\ref{sec:framework} is the conceptual intermediate: the pure cosine similarity
case corresponds to $\alpha = 0$ with $w_{\mathrm{char}} = w_{\mathrm{base}} \times \mathrm{sim}$.
In practice, $\alpha > 0$ is always used so that edges are never reduced to zero
by a low-similarity pairing alone.
The alpha parameter is tunable per plan configuration in the YAML config:
\begin{verbatim}
algorithm:
  weights: geographic
  character_signal: zone-membership
  character_alpha: 0.5
\end{verbatim}

\subsection{CLI Integration}

Each M-track signal is activated by the \texttt{--weights-override} flag with the
corresponding signal name:
\begin{verbatim}
bisect state --state NC --year 2020 \
  --weights-override zone-membership \
  --version v_zone_m6
\end{verbatim}
Multiple signals can be composed using the composite mode (M.8):
\begin{verbatim}
bisect state --state NC --year 2020 \
  --weights-override composite-community \
  --character-signals economic,zone-membership,housing \
  --version v_composite_m8
\end{verbatim}

\subsection{Interaction with Structure and Search Layers}

Community character weights interact with the structure layer in the following way:
the weights are computed once for the full state adjacency graph and remain constant
throughout all bisection rounds.
When the bisection tree recursively splits a subregion, METIS receives the subgraph
of that subregion together with its original edge weights; no recomputation of
character similarity is performed.
This means that the character weights influence not just the first bisection but
every subsequent level of the recursive partition.

Community character weights do not interact with the search layer (Layer 3).
All seed selection strategies---single, multi, convergence, percentile---operate on
the same weighted graph.
The choice of search strategy affects how many METIS trials are run and which result
is selected, but it does not alter the weights presented to METIS.

\subsection{Composability with County Weights}

The county weight mode (B.10) and the M.6 zone co-membership mode are related but
distinct.
County co-membership is a binary signal: two tracts are either in the same county
or not.
Zone co-membership (M.6) is a graded signal: tracts can share zero, one, two, three,
or all available zone types.
County subdivision membership (a dimension of M.6) is strictly more refined than
county membership in New England states, where the town is the primary administrative
unit and the county is a secondary one.

The two modes can be run independently or composed.
Running zone co-membership as the base mode subsumes the county signal when
county subdivision data is available, because county subdivision co-membership implies
county co-membership in every state.
In non-New England states where county is the primary unit, the county mode and the
M.6 county-subdivision dimension produce nearly identical results.
